\subsection{Sub-Baseline Scale}
\label{sec:sub-baseline}

Elements intentionally sized below the $6\,U$ text baseline --- such as checkboxes, radio buttons, range sliders, and progress bars --- follow a fixed proportional sub-scale:

\begin{equation}
\{2\,U,\; 4\,U,\; 6\,U\}.
\end{equation}

These values remain constant across density levels. The density factor $d$ does not affect sub-baseline elements because they have no wrapping-level padding to scale; their geometry is determined entirely by the base unit $U$.

The sub-scale provides three discrete heights:

\begin{center}
\begin{tabular}{cll}
\toprule
Height & px at $16\,\text{px}$ & Examples \\
\midrule
$2\,U$ & $8\,\text{px}$ & progress bar, popover arrow \\
$4\,U$ & $16\,\text{px}$ & checkbox, radio, range thumb, switch \\
$6\,U$ & $24\,\text{px}$ & text line (baseline) \\
\bottomrule
\end{tabular}
\end{center}

The boundary at $6\,U$ marks the transition: elements at or above this height enter the parametric model governed by Equation~\ref{eq:height}; elements below it are governed by the sub-scale.
